\sect{Приложение Б. Псевдокод и анализ On-line Time Warping}{appendix-oltw}
%
В настоящем приложении приведены подробности On-line Time Warping,
используемого как резервный трекер в~\Secref{realtime}. Алгоритм
основан на каузальной версии классической DTW-рекурсии
\eqref{eq:dtw-recursion}, в которой матрица накопленных стоимостей
$D$ заполняется не целиком, а столбец за столбцом по мере поступления
новых наблюдений.

Состояние OLTW описывается тремя величинами: текущей строкой
$i_{\text{cur}}$ (последняя обработанная позиция партитуры), текущим
столбцом $j_{\text{cur}}$ (число принятых наблюдений) и счётчиком
\textit{RunCount}, ограничивающим число последовательных шагов в одну
сторону. Параметр $R_{\max}$ устанавливает верхнюю границу для
RunCount; в нашей реализации $R_{\max}=3$.

\par\addvspace{2pt}
{\centering
\begin{minipage}{0.97\linewidth}
\footnotesize
\textbf{Алгоритм 2.} OLTW в каузальном режиме. \\[1pt]
\rule{\linewidth}{0.4pt}\vspace{1pt}\\
\textbf{Вход:} партитура $S=(p_1,\dots,p_N)$, ширина окна $W$,
ограничение $R_{\max}$. \\
1: $i_{\text{cur}}\gets 1,\,j_{\text{cur}}\gets 1$,
   \,RunCount $\gets 0$, \,direction $\gets 0$ \\
2: $D[1,1]\gets d(s_1,q_1)$ \\
3: \textbf{для} каждой поступающей ноты $q_j$ \textbf{выполнить} \\
4: \quad обновить столбец $j$ матрицы $D$ в окне
   $[\,i_{\text{cur}}-W,\,i_{\text{cur}}+W\,]$
   по~\eqref{eq:dtw-recursion} \\
5: \quad $(i^*,j^*)\gets\argminop\bigl\{D[i,j]:
   \;|\,i-i_{\text{cur}}|\le 1,\,|\,j-j_{\text{cur}}|\le 1\bigr\}$ \\
6: \quad \textbf{если} $(i^*,j^*)$ продолжает текущее
   direction \textbf{то} \\
7: \quad\quad RunCount $\gets$ RunCount $+\,1$ \\
8: \quad \textbf{иначе} \\
9: \quad\quad RunCount $\gets 1$;\,\, direction $\gets$
   шаг$(i_{\text{cur}}\!\to\!i^*,\,j_{\text{cur}}\!\to\!j^*)$ \\
10: \quad \textbf{конец если} \\
11: \quad \textbf{если} RunCount $\ge R_{\max}$ \textbf{то} \\
12: \quad\quad принудительно сменить direction
   на ортогональное, RunCount $\gets 0$ \\
13: \quad \textbf{конец если} \\
14: \quad $i_{\text{cur}}\gets i^*,\;j_{\text{cur}}\gets j^*$ \\
15: \quad \textbf{выдать} $\hat\varphi(j)\gets i_{\text{cur}}$ \\
16: \textbf{конец для} \\
\rule{\linewidth}{0.4pt}
\end{minipage}\par}
\addvspace{4pt}

Ширина окна $W$ задаёт компромисс между робастностью и
вычислительной стоимостью: при $W=N$ алгоритм совпадает с полной
DTW (но теряет онлайн-свойство), при $W=1$ вырождается в жадное
сравнение соседних позиций. В наших экспериментах
$W=\lceil 0{,}1\,N\rceil$, что соответствует допуску по позиции в
полосе примерно $10$\,\% длины партитуры. Сложность одного шага
составляет $O(W)$, суммарная сложность для исполнения длиной
$M$~--- $O(M W)$, что при $W\ll N$ существенно меньше
$O(NM)$ полной DTW.

Геометрический смысл механизма \textit{RunCount} проиллюстрирован
ниже. Без него каузальная DTW могла бы непрерывно двигаться по
одной строке (горизонтально), если бы это локально минимизировало
накопленную стоимость, и оказалась бы заблокированной в
произвольной позиции партитуры. Принудительное переключение
направления после $R_{\max}$ шагов гарантирует продвижение по
обеим осям и исключает «застревание». Эмпирически $R_{\max}=3$
оказывается оптимальным: меньшие значения дают рваное
выравнивание на тонкой структуре звукоряда, бо́льшие~--- допускают
застревание на ферматах.

В сочетании с HMM (см.~\Secref{realtime}) OLTW служит резервным
трекером, активируемым при низкой уверенности forward-распределения.
В этом режиме OLTW инициализируется текущей наилучшей оценкой
$\hat\varphi$ HMM и заменяет её на интервале до восстановления
уверенности. После восстановления HMM продолжается с
$\alpha$-вектором, переинициализированным в окне $\pm 4$ ноты
вокруг последнего значения OLTW.
